\documentclass[10pt,a4paper]{amsart}
\usepackage[margin=19mm]{geometry}
\usepackage{amsmath,amssymb,mathtools}
\usepackage[scr=rsfs,cal=euler]{mathalfa}
\usepackage{xfrac}
\usepackage[unicode,hidelinks,bookmarks=false]{hyperref}
\newcommand{\ie}{i.e.\ }
\newcommand{\cf}{cf.\ }
\newcommand{\resp}{respectively\ }
\newcommand{\Subsection}{\subsection}
\providecommand{\nobreakdash}{\nobreak\hbox{-}}
\setlength{\emergencystretch}{2em}
\multlinegap0pt
\usepackage[shortlabels]{enumitem}
\usepackage{xcolor}
\usepackage{amssymb}
\newtheorem{lemma}{Lemma}
\newcommand{\A}{\mathcal{A}}
\title{The first moment of quadratic Dirichlet $L$-functions in the even hyperelliptic ensemble}
\author{Hwanyup Jung}
\date{Source: arXiv:2609.12764v1 (CC BY 4.0)}
\begin{document}
\maketitle

\noindent\emph{Source and licence.} The first moment of quadratic Dirichlet $L$-functions in the even hyperelliptic ensemble. Version arXiv:2609.12764v1 (CC BY 4.0), distributed by arXiv under the Creative Commons Attribution 4.0 International licence, http://creativecommons.org/licenses/by/4.0/, as recorded in the arXiv licence metadata for this exact version and shown on the abstract page https://arxiv.org/abs/2609.12764v1. Full licence notice: \texttt{licenses/arxiv-2609.12764-CC-BY-4.0.txt}.\par\smallskip

\noindent\textit{Editorial scope.} Counted unit: the article's lemma lem:K-poles (source0.tex lines 1962--1968). The article's complete original proof of it is delivered with it (source0.tex lines 1970--2034, one \texttt{proof} environment of 65 lines). No mathematics is written, repaired or re-derived: the statement and the proof are the article's own lines, in the article's own notation, and no retained line is substituted or re-flowed.  Retained uncounted as the source-local context the proof needs: lines 1617--1698; lines 1671--1680; lines 1683--1693.  Nothing was omitted from any retained span, so the package carries no omission record.  No retained line contains a citation, so the wrapper carries no bibliography and no bibliography entry.  No retained line contains a figure environment, an image-inclusion command or drawing code, so no image is delivered and the package declares no asset dependency.  Editorial formatting only, in addition: an AMS \texttt{amsart} preamble that restates the article's own declarations and macros used by the retained lines, namely the environments \texttt{lemma} on separate counters, together with the standard packages those lines need.  The retained environments print in this document as Lemma 1 in order of appearance, whereas the article prints the same items with the numbering of its own full document.  The complete native source of the article as distributed in the arXiv e-print for this exact version is bundled unchanged as \texttt{sources/210148/source0.tex}.
\begin{lemma}[Euler-product factorization of $\mathcal B$]\label{lem:B}
For $q^{-2}<|z|<q^{-1}$ and $|w'|$ sufficiently small, the series
$\mathcal B(z,w')$ converges absolutely, and it admits the following
factorization.

\begin{enumerate}[label=\textup{(\arabic*)}]
\item
One has
\begin{equation}\label{eq:B-factorization}
\mathcal{B}(z,w')
=
\mathcal Z_{\A}(z)
\mathcal Z_{\A}(w')
\mathcal Z_{\A}(q(w')^2z)
\prod_{P\in\mathcal P}\mathcal{B}_P(z,w'),
\end{equation}
where, writing
$Z_P:=z^{d(P)}$ and $W_P:=(w')^{d(P)}$ for $P\in\mathcal P$,
the local factor is
\begin{equation}\label{eq:BP}
\mathcal B_P(z,w')
=
1+
\frac{
W_P\bigl(1-|P|^2Z_P^2\bigr)
-|P|(|P|-1)Z_PW_P^2
-|P|Z_P\bigl(1-|P|Z_P\bigr)W_P^3
}{
|P|^2Z_P-1
}.
\end{equation}
Moreover, the Euler product
$\prod_{P\in\mathcal P}\mathcal B_P(z,w')$
converges absolutely under the weaker restriction
\begin{equation}\label{eq:B-region}
|w'|<q|z|,
\qquad
|w'|<q^{-1/2},
\qquad
|w'z|<q^{-1}.
\end{equation}

\item
For each $P\in\mathcal P$, define
\begin{equation}\label{eq:DP}
\mathcal D_P(z,w')
:=
\mathcal B_P(z,w')
\left(
1-\frac{W_P}{|P|^2Z_P}
\right)
(1-W_P^2)^{-1}.
\end{equation}
Then one has
\begin{align}\label{eq:B-continuation}
\mathcal{B}(z,w')
&=
\mathcal Z_{\A}(z)
\mathcal Z_{\A}(w')
\mathcal Z_{\A}(q(w')^2z)
\mathcal Z_{\A}\!\left(\frac{w'}{q^2z}\right)
\mathcal Z_{\A}((w')^2)^{-1}
\prod_{P\in\mathcal P}\mathcal{D}_P(z,w').
\end{align}
The right-hand side furnishes a meromorphic continuation of $\mathcal{B}(z,w')$
to the region
\begin{equation}\label{eq:D-region}
|z|>q^{-2},
\qquad
|w'|^2<q|z|,
\qquad
|w'|<q^3|z|^2,
\qquad
|w'|<1,
\qquad
|w'z|<q^{-1},
\end{equation}
where the Euler product
$\prod_{P\in\mathcal P}\mathcal{D}_P(z,w')$
converges absolutely and uniformly on compact subsets.
\end{enumerate}
\end{lemma}

\begin{align}
\mathcal{B}(z,w')
&=
\mathcal Z_{\A}(z)
\mathcal Z_{\A}(w')
\mathcal Z_{\A}(q(w')^2z)
\mathcal Z_{\A}\!\left(\frac{w'}{q^2z}\right)
\mathcal Z_{\A}((w')^2)^{-1}
\prod_{P\in\mathcal P}\mathcal{D}_P(z,w').
\end{align}

\begin{equation}
|z|>q^{-2},
\qquad
|w'|^2<q|z|,
\qquad
|w'|<q^3|z|^2,
\qquad
|w'|<1,
\qquad
|w'z|<q^{-1},
\end{equation}

\begin{lemma}[Poles of $K_{N,\bullet}$]
\label{lem:K-poles}
For $|z|=q^{-3/2}$, $K_{N,\bullet}(z,\cdot)$ is meromorphic on the disc
$|w|<1$, with no poles there other than the five points $w=0$,
$w=\pm z^{1/2}$ of modulus $q^{-3/4}$, and $w=\pm q^{2}z^{3/2}$ of modulus
$q^{-1/4}$.
\end{lemma}

\begin{proof}
On $|z|=q^{-3/2}$ the substitution $w'=w/(qz^{1/2})$ gives
$|w'|=q^{-1/4}|w|$, and for $|w|<1$ the pair $(z,w')$ satisfies the five
conditions of \eqref{eq:D-region}, the binding one being $|w'|^2<q|z|$,
which reads $|w|<1$; these conditions involve $w'$ only through $|w'|$, so
they hold for $-w'$ as well, hence for both terms of
$\mathcal B^{\bullet}(z,w')$. By Lemma~\ref{lem:B}(2) the right-hand side of
\eqref{eq:B-continuation} therefore represents $\mathcal B(z,\pm w')$ on
$|w|<1$ as a meromorphic function whose only poles are those of the explicit
zeta factors. As the remaining factors of $K_{N,\bullet}$ are rational in
$w$, it follows that $K_{N,\bullet}(z,\cdot)$ is meromorphic on $|w|<1$.

It remains to locate its poles. Inserting \eqref{eq:B-continuation} into
$\mathcal B^{\bullet}$ and undoing the substitution $w=qz^{1/2}w'$ displays
every $w$-dependent factor of $K_{N,\bullet}$ at once: with
$\varrho\in\{\pm1\}$ indexing the two terms of $\mathcal B^{\bullet}$,
$$
K_{N,\bullet}(z,w)
=
\frac{z^{-\nu_\bullet/2}\,\mathcal Z_{\A}(z)
\left(1-\frac{w^2}{qz}\right)}
{2\,(1-w^2)(1-w)(1-q^{-1/2}w)\,w^{N+1}}
\sum_{\varrho\in\{\pm1\}}
\frac{\varrho^{\nu_\bullet}
\prod_{P\in\mathcal P}\mathcal D_P(z,\varrho w')}
{\left(1-\frac{\varrho w}{z^{1/2}}\right)
\left(1-\frac{\varrho w}{q^2z^{3/2}}\right)}.
$$
Here the four zeta factors of \eqref{eq:B-continuation} that involve $w$
have been evaluated as follows, the last column recording the moduli on
$|z|=q^{-3/2}$:
\begin{center}
\renewcommand{\arraystretch}{1.4}
\begin{tabular}{c c c c}
factor of \eqref{eq:B-continuation} & in the variable $w$ & singular at &
$|w|$\\
\hline
$\mathcal Z_{\A}(\varrho w')$
& $\left(1-\dfrac{\varrho w}{z^{1/2}}\right)^{-1}$
& $w=\varrho z^{1/2}$
& $q^{-3/4}$\\

$\mathcal Z_{\A}\!\left(\dfrac{\varrho w'}{q^2z}\right)$
& $\left(1-\dfrac{\varrho w}{q^2z^{3/2}}\right)^{-1}$
& $w=\varrho q^2z^{3/2}$
& $q^{-1/4}$\\

$\mathcal Z_{\A}\bigl(q(w')^2z\bigr)$
& $\left(1-w^2\right)^{-1}$
& $w=\pm1$
& $1$\\

$\mathcal Z_{\A}\bigl((w')^2\bigr)^{-1}$
& $1-\dfrac{w^2}{qz}$
& nowhere
& \\
\end{tabular}
\end{center}
The Euler product $\prod_P\mathcal D_P(z,\varrho w')$ is holomorphic on
$|w|<1$ by Lemma~\ref{lem:B}(2), and the remaining factors $(1-w)^{-1}$,
$(1-q^{-1/2}w)^{-1}$ and $w^{-(N+1)}$ of $K_{N,\bullet}$ are singular only
at $w=1$, $w=q^{1/2}$ and $w=0$, respectively. Of all these singularities
those at $w=\pm1$ and $w=q^{1/2}$ have modulus $\ge1$, so the only poles in
$|w|<1$ are $w=0$, $w=\pm z^{1/2}$ and $w=\pm q^{2}z^{3/2}$.
\end{proof}

\end{document}
